\documentclass[11pt]{article}
\usepackage[margin=1in]{geometry}
\usepackage{booktabs}
\usepackage{longtable}
\usepackage{hyperref}
\usepackage{xcolor}
\usepackage{amsmath}
\usepackage{parskip}

\title{Independent Replication Report: \\
\emph{Arthrobacter} sp.\ PAMC25564 Complete Genome and CAZyme-Based Cold Adaptation \\
{\large Han et al.\ 2021, \textit{BMC Genomics} 22:403 (PMID 34078272)}}
\author{BVBRC-59 Replication Track --- Ollie (rick-stevens-ai)}
\date{2026-07-05}

\begin{document}
\maketitle

\begin{abstract}
We independently rerun the core computational claims of Han \emph{et al.} (2021) on the Antarctic cryoconite isolate \emph{Arthrobacter} sp.\ PAMC25564 using only public NCBI data and open-source tooling. All six primary genome-summary statistics (length, GC, total genes, CDS, pseudogenes, rRNA, tRNA) reproduce \emph{exactly} against the paper-contemporaneous GenBank annotation. An independent HMMER/dbCAN-HMMdb-V9 CAZyme rerun yields 102 vs.\ the paper's 108, with all glycogen/trehalose cold-adaptation signature families (GH1, GH13\_11, GH13\_26, GH65, GH77, CBM48) fully recovered. An independent LLM judge (Argo gpt-5.2) returns \textbf{REPLICATED} with coverage 9/10 and agreement 8/10.
\end{abstract}

\section{Verdict}
\textbf{REPLICATED} \quad (coverage 9/10, agreement 8/10, judge: Argo gpt-5.2).

\section{Scope}
Complete genome of \emph{Arthrobacter} sp.\ PAMC25564 plus comparative genomics across 16--26 \emph{Arthrobacter}/\emph{Pseudarthrobacter} complete genomes, emphasising: (a) genome characterisation, (b) the CAZyme complement classified by dbCAN2, and (c) the glycogen/trehalose metabolic CAZyme families proposed to underpin cold adaptation.

\section{Data}
\begin{itemize}
  \item \textbf{Focal genome:} CP039290.1 (RefSeq NZ\_CP039290.1); assembly GCA\_004798705.1 / GCF\_004798705.1 (ASM479870v1); BioProject PRJNA531357; BioSample SAMN11356967. PacBio Sequel, HGAP v4. Submitter: Korea Polar Research Institute.
  \item \textbf{Focal proteome:} GenBank PGAP proteins, 3{,}613 sequences.
  \item \textbf{Comparators (verified real/public, sampled):} \emph{Arthrobacter} sp.\ 24S4-2 CP040018.1, PAMC25486 CP007595.1, ZXY-2 CP017421.1; \emph{Crystallibacter/Arthrobacter crystallopoietes} DSM 20117 CP018863.1; \emph{Pseudarthrobacter phenanthrenivorans} Sphe3 CP002379.1.
  \item \textbf{dbCAN HMM DB:} dbCAN-HMMdb-V9 (pro.unl.edu; HMMER3), substituted for the paper's dbCAN2/V8-era DB (bcb.unl.edu offline post-cyberattack).
\end{itemize}

\section{Methods}
\begin{longtable}{@{}clll@{}}
\toprule
\# & Step & Paper & This rerun \\
\midrule
M1 & Genome length + GC        & assembly announcement       & Parsed CP039290.1 FASTA (Python) \\
M2 & Gene / CDS / RNA counts   & NCBI PGAP annotation        & NCBI Datasets v2 (original GenBank 2019-04-11) + efetch \\
M3 & Proteome                  & PGAP CDS                    & NCBI Datasets PROT\_FASTA (3{,}613 proteins) \\
M4 & CAZyme classification     & dbCAN2 (HMMER + Hotpep + DIAMOND) & HMMER 3.4 vs dbCAN-HMMdb-V9 + canonical parser \\
M5 & Cold-adapt families       & dbCAN2 + manual (Table 2)   & Family membership from M4 domtbl \\
M6 & Comparator availability   & 16--26 genomes              & NCBI esummary of sampled accessions \\
\bottomrule
\end{longtable}

Tooling: HMMER 3.4 (conda \texttt{antismash} env, uicgpu01, 16 CPU); Python 3.8; NCBI E-utilities + Datasets v2 REST; Europe PMC OA XML. All endpoints free. LLM judge = Argo gpt-5.2 (localhost:44497). CAZyme parser: canonical dbCAN \texttt{hmmscan-parser} filter ($E<10^{-15}$ if alignment $>80$ aa else $E<10^{-5}$; HMM coverage $>0.35$; overlap $>0.5$ resolution).

\section{Claims Table}
\begin{longtable}{@{}clllll@{}}
\toprule
ID & Claim & Type & Testable & Tested & Result \\
\midrule
C1  & Circular 4{,}170{,}970 bp                & quantitative & yes & yes & \textbf{EXACT} \\
C2  & GC 66.74\%                                & quantitative & yes & yes & MATCH ($\Delta$0.03) \\
C3a & 3{,}829 total genes                       & quantitative & yes & yes & \textbf{EXACT} \\
C3b & 3{,}613 CDS                               & quantitative & yes & yes & \textbf{EXACT} \\
C3c & 147 pseudogenes                           & quantitative & yes & yes & \textbf{EXACT} \\
C3d & 15 rRNA                                   & quantitative & yes & yes & \textbf{EXACT} \\
C3e & 51 tRNA                                   & quantitative & yes & yes & \textbf{EXACT} \\
C4  & 16--26 complete \emph{Arthrobacter} set    & data availability & yes & sampled & VERIFIED \\
C5  & 108 CAZymes (33/45/23/5/2/0)              & quantitative & yes & yes & CLOSE (102: 34/43/16/5/9/0) \\
C6  & GH1/GH13\_11/GH13\_26/GH65/GH77/CBM48     & qualitative  & yes & yes & \textbf{FULL MATCH} \\
\bottomrule
\end{longtable}

\section{Results}

\subsection{Genome statistics (M1--M3)}
\begin{tabular}{@{}lccc@{}}
\toprule
Metric & Paper & Rerun & Verdict \\
\midrule
Length (bp)    & 4{,}170{,}970 & 4{,}170{,}970 & EXACT \\
GC (\%)        & 66.74         & 66.71         & $\Delta$0.03 \\
Total genes    & 3{,}829       & 3{,}829       & EXACT \\
Protein-coding & 3{,}613       & 3{,}613       & EXACT \\
Pseudogenes    & 147           & 147           & EXACT \\
rRNA           & 15            & 15            & EXACT \\
tRNA           & 51            & 51            & EXACT \\
\bottomrule
\end{tabular}

\medskip
\emph{Annotation-version note.} The current RefSeq re-annotation (RS\_2024\_05\_22) has drifted (3{,}863 / 3{,}718 / 75). Using the paper-contemporaneous \textbf{original GenBank annotation (2019-04-11)} reproduces all six counts exactly --- confirming annotation \emph{version} as the sole source of any discrepancy.

\subsection{CAZymes (M4--M5)}
\begin{tabular}{@{}lcc@{}}
\toprule
Class & Paper (dbCAN2) & Rerun (HMMER/dbCAN-HMMdb-V9) \\
\midrule
Total & 108 & 102 \\
GH    & 33  & 34 \\
GT    & 45  & 43 \\
CE    & 23  & 16 \\
AA    & 5   & 5  \\
CBM   & 2   & 9  \\
PL    & 0   & 0  \\
\bottomrule
\end{tabular}

\medskip
GH/GT/AA/PL match within 0--2; total within 6. CE (lower) and CBM (higher) shifts are the expected consequence of (i) dbCAN DB version V9 vs V8 and (ii) HMMER-only rerun here vs dbCAN2's 3-signature overview in the paper. CBM and CE calls are the most version-sensitive.

\textbf{Cold-adaptation signature families (paper Table 2) --- independently recovered:} GH1 ($\beta$-glucosidase); GH13 with 7 subfamilies including GH13\_11 (glycogen debranching) and GH13\_26; GH65 ($\alpha$-trehalose phosphorylase); GH77 (4-$\alpha$-glucanotransferase); CBM48. Fully consistent with the paper's glycogen + trehalose metabolism conclusion.

\subsection{Comparative dataset (M6)}
Sampled comparator accessions (CP040018.1, CP007595.1, CP017421.1, CP018863.1, CP002379.1) all resolve to real, complete public genomes of the named strains, confirming the comparative dataset the paper's conclusions rest on is fully available.

\section{BV-BRC Mapping}
The paper's workflow maps cleanly onto \textbf{BV-BRC Comprehensive Genome Analysis} (assembly stats + PGAP-equivalent annotation $\rightarrow$ genome length/GC/gene/RNA counts) and \textbf{BV-BRC Specialty Genes / protein-family services} (CAZyme classification). The rerun executes the equivalent steps with open tooling on the identical public genome.

\section{Genuine Critique}
This section records the honest weaknesses of both the paper and the rerun, and where a reader should be cautious.

\subsection{Weaknesses of the rerun}
\begin{itemize}
  \item \textbf{Tool substitution is not tool equivalence.} We rerun CAZyme classification with HMMER-only against dbCAN-HMMdb-V9, but the paper used the full dbCAN2 3-signature pipeline (HMMER + Hotpep + DIAMOND) against a V8-era DB. HMMER-only is the most conservative leg of that pipeline; the observed CBM \emph{increase} (2$\rightarrow$9) and CE \emph{decrease} (23$\rightarrow$16) are consistent with (i) different filtering thresholds and (ii) V9's expanded CBM HMM set. We assert this is version-bound noise, not disagreement, but we cannot fully separate the two effects without also running Hotpep+DIAMOND on a V8 snapshot --- which is currently unavailable due to the bcb.unl.edu cyberattack outage.
  \item \textbf{Comparative pan-CAZyme analysis not re-run end-to-end.} We verified the comparator dataset is real and reachable via NCBI, but we did not reproduce the pan-CAZyme comparative counts (e.g.\ the paper's comparative figures). The paper's biological conclusion --- that PAMC25564's glycogen/trehalose CAZyme enrichment is distinctive among mesophilic \emph{Arthrobacter} --- is therefore \emph{consistent with} but not \emph{independently retested by} this replication.
  \item \textbf{No cold-shock phenotype work.} The paper's cold-adaptation claim is inferred from CAZyme complement; neither the paper nor the rerun ties specific CAZymes to measured cold-growth phenotypes. This is an inherent limitation of the underlying study design and carries over.
  \item \textbf{Annotation-drift caveat.} We reproduced counts using the 2019-04-11 GenBank annotation; anyone repeating this in the future against the current RefSeq annotation will get different numbers (3{,}863 / 3{,}718 / 75) and may falsely conclude the paper is wrong. This is a version-pinning issue, not a claim error, but it is a real replication hazard.
\end{itemize}

\subsection{Weaknesses of the paper itself (honest read)}
\begin{itemize}
  \item The cold-adaptation claim is largely \emph{genomic association} (family presence + counts) rather than \emph{functional demonstration}. Even a full CAZyme catalogue does not, on its own, prove those enzymes are engaged during cold growth. Transcriptomics or enzyme-activity assays would be required to elevate the claim from correlational to mechanistic.
  \item The comparative set (16--26 genomes) is small relative to the current diversity of sequenced \emph{Arthrobacter}; conclusions about "distinctive" CAZyme complements could shift as more Antarctic isolates are sequenced.
  \item The paper's dbCAN2 version and parameter set are not fully machine-reproducible from the manuscript text (typical for the era); a reader must accept published counts on trust or re-derive them.
\end{itemize}

\subsection{What this replication does and does not prove}
\emph{Proves:} the assembly and annotation the paper reports are recoverable exactly from public archives against the correct-vintage annotation; the CAZyme complement (approximately) and the signature glycogen/trehalose families are independently recoverable; the comparator dataset is real.

\emph{Does not prove:} that PAMC25564's CAZyme complement is causally responsible for its cold adaptation; that the paper's dbCAN2/V8 numbers would be reproduced \emph{exactly} using the same tool at the same version.

\section{Limitations}
\begin{itemize}
  \item CAZyme rerun used HMMER-only vs dbCAN-HMMdb-V9 rather than the full dbCAN2/V8 3-tool overview; totals within 6, all biologically-relevant families recovered, but exact per-class parity (esp.\ CE/CBM) is version-bound.
  \item Comparator set verified for availability (sampled) rather than fully re-run through pan-CAZyme comparison.
\end{itemize}

\section{Verdict}
\textbf{Verdict: REPLICATED.} All six primary genome statistics reproduced \emph{exactly} on public NCBI data against the paper-contemporaneous annotation; independent HMMER/dbCAN rerun gave 102 CAZymes vs.\ the paper's 108 with all glycogen/trehalose cold-adaptation families (GH1, GH13\_11, GH13\_26, GH65, GH77, CBM48) recovered; LLM judge REPLICATED, coverage 9/10, agreement 8/10.

\end{document}
